\subsection{Functions of positive type on a locally compact commutative group}

In this number, G is a locally compact commutative group and \(\mu\) denotes a Haar measure on G. We denote by \(\widehat{G}\) the dual group of G and by \(\widehat{\mu}\) the dual Haar measure of \(\mu\) (def. 4 of II, p. 214). We denote by F the Fourier transform on the Banach space \(\mathcal{M}^{1}(G)\) of bounded measures on G.

Since \(\mathcal{C}_0(\mathrm{G})\) is the closure of \(\mathcal{K}(\mathrm{G})\) in \(\mathcal{C}_b(\mathrm{G})\) , the Banach space \(\mathcal{M}^1 (\mathrm{G})\) , dual of \(\mathcal{K}(\mathrm{G})\) endowed with the topology of \(\mathcal{C}_b(\mathrm{G})\) , is identified with the dual of \(\mathcal{C}_0(\mathrm{G})\) (cf. INT, III, p. 56, § 1, no. 8). We endow \(\mathcal{M}^1 (\mathrm{G})\) with the weak topology associated with this duality.

Lemma 9. --- The Fourier transform is a continuous mapping from \(\mathcal{M}^{1}(\mathrm{G})\) equipped with the weak topology to \(\mathcal{C}_{b}(\widehat{\mathrm{G}})\) endowed with the topology induced by the weak topology \(\sigma(\mathrm{L}^{\infty}(\widehat{\mathrm{G}}),\mathrm{L}^{1}(\widehat{\mathrm{G}}))\) .

Let \(\mathfrak{F}\) be a filter on \(\mathcal{M}^1 (\mathrm{G})\) weakly converging to a bounded measure \(\nu\) on G. For every \(\varphi \in \mathrm{L}^1 (\widehat{\mathrm{G}})\) , we have \(\mathcal{F}(\varphi)\in \mathcal{C}_0(\mathrm{G})\) (prop. 4 of II, p. 209), hence

\[
\begin{array}{r l} & {\int_ {\widehat {\mathrm{G}}} \varphi \mathcal {F} (\nu) d \widehat {\mu} = \int_ {\mathrm{G}} \mathcal {F} (\varphi) d \nu} \\ & {\qquad = \lim _ {\eta , \mathfrak {F}} \int_ {\mathrm{G}} \mathcal {F} (\varphi) d \eta = \lim _ {\eta , \mathfrak {F}} \int_ {\widehat {\mathrm{G}}} \varphi \mathcal {F} (\eta) d \widehat {\mu},} \end{array}
\]

by the transposition property of the Fourier transform (prop. 13 of II, p. 221). The lemma follows.

THEOREM 5 (Bochner). --- A continuous function \(\varphi\) on \(\widehat{\mathbf{G}}\) belongs to \(\operatorname{Pos}(\widehat{\mathbf{G}})\) if and only if there exists a bounded positive measure \(\nu\) on \(\mathbf{G}\) such that \(\varphi = \mathcal{F}(\nu)\) .

Let \(\nu \in \mathcal{M}^{1}(\mathrm{G})\) be a positive measure. Its Fourier transform is continuous. For every finite family \((x_{i})_{i\in \mathrm{I}}\) in \(\widehat{\mathbf{G}}\) and every finite family \((t_i)_{i\in \mathrm{I}}\) of complex numbers, it follows that

\[
\begin{array}{r} \sum_ {i \in \mathrm{I}} \sum_ {j \in \mathrm{I}} \bar {t} _ {i} t _ {j} \mathcal {F} (\nu) (x _ {i} ^ {- 1} x _ {j}) = \sum_ {i \in \mathrm{I}} \sum_ {j \in \mathrm{I}} \bar {t} _ {i} t _ {j} \int_ {\mathrm{G}} x _ {i} \overline {{x}} _ {j} d \nu \\ = \int_ {\mathrm{G}} \left| \sum_ {i \in \mathrm{I}} \bar {t} _ {i} x _ {i} \right| ^ {2} d \nu \geqslant 0, \end{array}
\]

since \(\nu\) is a positive measure. The Fourier transform of \(\nu\) is therefore a function of positive type on \(\hat{G}\) (th. 1 of V, p. 432, (ii)).

Let us prove the converse. Endow the set \(\operatorname{Pos}_{\leqslant1}(\widehat{\mathbf{G}})\) with the weak topology, induced as before by the weak topology \(\sigma(\mathrm{L}^{\infty}(\widehat{\mathrm{G}}),\mathrm{L}^{1}(\widehat{\mathrm{G}}))\) . The set \(\operatorname{Pos}_{\leqslant1}(\widehat{\mathbf{G}})\) is compact and convex (cor. of prop. 13 of V, p. 448). By prop. 14 of V, p. 450, prop. 11 of V, p. 444, and cor. 7 of V, p. 390, the extreme points of \(\operatorname{Pos}_{\leqslant1}(\widehat{\mathbf{G}})\) are the zero function and the elements of \(\widehat{G}\) .

Let N be the set of positive bounded measures of mass \(\leqslant 1\) on G; it is compact in \(\mathcal{M}^{1}(G)\) endowed with the weak topology (EVT, III, p. 17, cor. 3). By Lemma 9 and the first part of the proof, the Fourier transform on G defines, by passing to subspaces, a continuous map from \(\mathcal{N}\) to \(\mathrm{Pos}_{\leqslant 1}(\widehat{\mathrm{G}})\) ; by homogeneity, it suffices to show that this map is surjective. The image \(\mathcal{F}(\mathcal{N})\) of \(\mathcal{N}\) by the Fourier transform is convex and compact; it contains the zero function and the elements of \(\widehat{\mathrm{G}}\) (indeed, these are of the form \(\mathrm{ev}_x\colon \chi \mapsto \chi (x)\) for an element \(x\) of G by th. 2 of II, p. 220, and we have \(\mathrm{ev}_x = \mathcal{F}(\varepsilon_{x^{-1}})\) ). The set \(\mathcal{F}(\mathcal{N})\) therefore contains the extreme points of \(\mathrm{Pos}_{\leqslant 1}(\mathrm{G})\) , whence \(\mathcal{F}(\mathcal{N}) = \mathrm{Pos}_{\leqslant 1}(\mathrm{G})\) by the Krein-Milman theorem (EVT, II, p. 59, th. 1). This concludes the proof.

When \(G = R^{k}\) for an integer \(k \in N\) , this theorem corresponds to prop. 11 of INT, IX, p. 94, § 6, no. 12.

